%%
%% fall-lecture01.tex
%% 
%% Made by Alex Nelson <pqnelson@gmail.com>
%% Login   <alex@lisp>
%% 
%% Started on  2026-10-13T11:04:01-0700
%% Last update 2026-10-13T11:04:01-0700
%% 

\lecture{}

\begin{node}
This quarter will focus on topological aspects of characteristic
classes. Next quarter will focus on geometric aspects of
characteristic classes. The textbook we will loosely follow is Milnor
and Stasheff~\cite{milnor1974characteristic}.
\end{node}

\begin{node}
The easiest way to construct a space out of old spaces $B$ and $Y$ is
to use their Cartesian product $B\times Y$. For example, the torus is
$T^{2}=S^{1}\times S^{1}$. But not all spaces can be written as the
trivial product---the M\"{o}bius band is one counter-example. But we
can \emph{informally} write the M\"{o}bius strip as the \emph{twisted}
interval bundle $S^{1}\widetilde{\times}I$ over $S^{1}$ (the base
space is still $S^{1}$ but we now have some twisting). It is locally
trivial $p^{-1}U\iso U\times I$ for sufficiently small open sets
$U$. This is generalized by the notion of fiber bundles.
\begin{enumerate}
\item We have a projection map $\pi\colon E\to B$
\item Locally trivializable: we can write $B=\bigcup_{k}U_{k}$ using a
  local cover, there exists a homeomorphism
  $\varphi_{k}\colon\pi^{-1}U_{k}\to U_{k}\times Y$ where $Y$ is the fiber
\item Gluings are given by the transition maps for any $U_{i}\cap U_{j}\neq\emptyset$,
  writing $U_{ij}=U_{i}\cap U_{j}$, then
\begin{equation}
\vcenter{\xymatrix{ & \ar[dl]^{\varphi_{i}}\pi^{-1}(U_{ij})\ar[dr]^{\varphi_{j}} & \\
U_{ij}\times Y\subset U_{i}\times Y\ar[rr]^{\varphi_{j}\circ\varphi_{i}^{-1}}& &
U_{ij}\times Y\subset U_{j}\times Y}}
\end{equation}
is a fiber-wise homeomorphism where $\varphi_{j}\circ\varphi_{i}^{-1}|_{U_{ij}}=\id$.
\end{enumerate}
We usually look at $Y=\RR^{n}$ or $Y=\CC^{n}$ for real (resp.,
complex) fiber bundles.
\end{node}

\begin{node}
An important class of examples: Let $M$ be a smooth $n$-manifold, then
the tangent bundle $\T M$ is $2n$-dimensional. For example, $\T S^{2}$,
there is a natural map
\begin{equation}
\pi\colon TS^{2}\to S^{2}
\end{equation}
sending $(x,v)\mapsto x$. This gives us a vector bundle.
\end{node}

\begin{node}
There are many different ways to study vector bundles. We usually
write it like
\begin{equation}
\vcenter{\xymatrix{Y\ar[r]& E\ar[d]\\
& B}}
\end{equation}
We can think of it as a continuous family of vector spaces
parametrized by the base space $B$. Then the operations for the direct
sum $Y_{1}\oplus Y_{2}$, tensor product $Y_{1}\otimes Y_{2}$,
$\hom(Y_{1},Y_{2})$, dual space $Y^{*}$, and linear functions
$\extp^{k}Y$ on vector spaces can be extended to  vector bundles.
\end{node}

\begin{node}
Studying vector bundles can give you some information about the base
manifold. This perspective will be implicit throughout the entire class.
\end{node}

\subsection{Characteristic classes}

\begin{node}
The main characteristic classes are the Stiefel--Whitney class, the
Pontryogin class, and the Chern class. These are the most important ones.
\end{node}

\begin{node}
Consider a real rank-$r$ bundle $E$ over $B$. Then the
Stiefel--Whiteney classes are
\begin{equation}
w_{i}(\T B)\in H^{i}(B,\ZZ_{2})
\end{equation}
for $i=0,\dots,r$. The $\ZZ_{2}$ coefficients means we can ignore
orientation. 
\end{node}

\begin{node}
The Pontryogin class $p_{i}\in H^{4i}(B,\ZZ)$ where $i=0,\dots,[r/2]$.
\end{node}

\begin{node}
The Chern class needs a complex structure, so let $V$ be a complex
rank-$r$ bundle over $B$ (so the fiber is $\CC^{r}$), then the Chern
class is
\begin{equation*}
C_{i}\in H^{2i}(B,\ZZ),
\end{equation*}
for $i=0,\dots,r$.
\end{node}

\begin{example}
The top Stiefel--Whitney class. Consider the tangent bundle $\T M$,
consider $w_{n}(\T M)\in H^{n}(M;\ZZ_{2})$. We know we have
$[M]_{2}\in H_{n}(M;\ZZ_{2})$ mod-2 fundamental class which is
nonzero. We have
\begin{equation}
\langle w_{n}(\T M),[M]_{2}\rangle=\chi(M)\bmod{2}
\end{equation}
is the Euler class module 2. For example, $M=\RP^{2}$, so
$[M]_{2}=[\RP^{2}]_{2}$, so $\chi(\RP^{2})=1$. Then $w_{2}(\RP^{2})\neq0$.
\end{example}

\begin{example}
Let $E$ be a holomorphic $r$-bundle over $M$ a closed oriented surface
of genus $g$ (``oriented'' may not be needed?). Then we have the first
Chern class $c_{1}\in H^{2}(M,\ZZ)$. The important conclusion (the
Hurewicz--Riemann--Roch theorem) is that
\begin{equation}
\dim_{\CC}H^{0}(M,E)-\dim_{\CC}H^{1}(M,E)=\langle c_{1}[M]\rangle+r(1-g),
\end{equation}
where $H^{0}(M,E)$ describes the space of holomorphic sections,
$H^{1}(M,E)$ is the sheaf cohomology (viz., a nonzero integer
dimension). Then
\begin{equation}
\dim_{\CC}H^{0}(M,E)\geq\langle c_{1},[M]\rangle+r(1-g).
\end{equation}
So this gives you information about this bundle (viz., the existence
of holomorphic sections).
\end{example}

\begin{node}[Classifying spaces]
We want to classify all rank-$r$ vector bundles over a fixed base
space $B$. If $X$ is a CW complex, we know we have a bijection
\begin{equation}
H^{i}(X,\ZZ)\stackrel{1:1}{\longleftrightarrow}[X,K(\ZZ,i)]
\end{equation}
where $K(\ZZ,i)$ is an Eilenberg--MacLane space, and $[X,Y]$ is the
set of homotopy equivalence classes of maps from $X$ to $Y$. Then we
want some analogous space replacing $K(\ZZ,i)$ for vector bundles,
which will turn out to be
\begin{equation}
\left\{\begin{matrix}
\mbox{isomorphism classes of}\\
\mbox{real rank-$k$ vector}\\
\mbox{bundles over $X$}\end{matrix}\right\}\stackrel{1:1}{\longleftrightarrow}[X,G_{k}(\RR)]
\end{equation}
where $G_{k}(\RR)=\bigcup_{n=k}G_{k}(\RR^{n})$ is the
infinite-dimensional Grassmannian manifold. For $k=1$, we see
$G_{1}(\RR)=\bigcup_{N=1}G_{1}(\RR^{N})\iso\RP^{\infty}$.
\end{node}

\begin{node}[Chern--Weil Theory]
We will try to represent characteristic classes using curvature of
vector bundle. For example, the Gauss--Bonnet theorem: a closed
oriented surface $M$ has
\begin{equation}
\int_{M}\kappa=2\pi\chi(M)
\end{equation}
where $\kappa$ is the Gaussian curvature, and $\chi(M)$ is the Euler
characteristic of $M$. This can be generalized to the
Chern--Gauss--Bonnet's theorem
\begin{equation}
\chi(M) = \frac{1}{(2\pi)^{n}}\int_{M}\Pfaff(\Omega)
\end{equation}
where $\Pfaff(\Omega)$ is the Pfaffian of $\Omega$, and $\Omega$ is
the curvature 2-form, $M$ is an $2n$-manifold. Then $\Omega$ can be
written with its components $\Omega_{ij}$ which are 2-forms. Here
\begin{equation}
\Omega_{ij}(X,Y)=\langle R(X,Y)e_{j},e_{i}\rangle.
\end{equation}
For $n=1$, $\Omega$ is a $2\times 2$ anti-symmetric matrix and
$\Pfaff(\Omega)=\Omega_{12}=\kappa\,\D A$ is the Gaussian curvature
multiplied by the area 2-form $\D A$. Then we recover the
Gauss--Bonnet theorem.

For $n=4$, we can compute
\begin{equation}
\Pfaff(\Omega)=\Omega_{12}\wedge\Omega_{34}-\Omega_{13}\wedge\Omega_{24}+\Omega_{14}\wedge\Omega_{23}.
\end{equation}
If $M$ has constant sectional curvature $K$, then we find
\begin{equation}
\Pfaff(\Omega)=3K^{2}\,\D A,
\end{equation}
and
\begin{equation}
\chi(M)=\frac{1}{4\pi^{2}}\int_{M}3 K^{2}\,\D A.
\end{equation}
For $\sphere{4}$, $\chi(\sphere{4})=2$ and we can conclude
\begin{equation}
\frac{4\pi^{2}}{3}2=\int K^{2}\,\D A.
\end{equation}
\end{node}